% Generated draft: requires mathematical and source review.
\documentclass[11pt]{article}
\usepackage[margin=25mm]{geometry}
\usepackage{fontspec}
\setmainfont{Latin Modern Roman}
\usepackage{amsmath,amssymb,mathtools,unicode-math}
\setmathfont{Latin Modern Math}
\usepackage{enumitem,needspace,xurl,hyperref}
\setlength{\parindent}{0pt}
\setlength{\parskip}{5pt}
\setlength{\emergencystretch}{3em}
\newcommand{\N}{\mathbb{N}}
\newcommand{\Z}{\mathbb{Z}}
\newcommand{\Q}{\mathbb{Q}}
\newcommand{\R}{\mathbb{R}}
\newcommand{\C}{\mathbb{C}}
\newcommand{\F}{\mathbb{F}}
\newcommand{\A}{\mathbb{A}}
\newcommand{\PP}{\mathbb{P}}
\newcommand{\E}{\mathbb{E}}
\newcommand{\eps}{\varepsilon}
\newcommand{\dd}{\,\mathrm{d}}
\begin{document}
\section*{UnsolvedMath ID: 700882}
\textbf{Source problem code:} catalogue entry
\subsection{Definitions and mathematical statement}
For integers \(k\geq 2\) and \(n\geq 1\), define the following assertion \(P(k,n)\). For a finite group \(G\), let \(1_G\) denote its identity and let
\[
\exp(G)=\operatorname{lcm}\{\operatorname{ord}(g):g\in G\}.
\]
Then \(P(k,n)\) is the statement
\[
\begin{gathered}
\exists N\in\Z_{\geq 1}\ \forall\text{ finite groups }G\ \forall(a_1,\ldots,a_k)\in G^k,\\
\left[\ G=\langle a_1,\ldots,a_k\rangle\ \land\
\bigl(\forall(x_1,\ldots,x_N)\in
(\{a_1,\ldots,a_k\}\cup\{1_G\})^N,\
(x_1\cdots x_N)^n=1_G\bigr)\ \right]
\Longrightarrow \exp(G)\leq n.
\end{gathered}
\]
Here \(N=N(k,n)\) must be independent of the finite group and its generating family. The factors in the tested products are generators or the identity; inverses are not separately included. Identity factors allow every product of at most \(N\) generators to be tested. No minimality requirement on the generating family is stated.

The source asks whether \(P(k,n)\) holds for every \(k\geq 2\) and every ``(prime) number'' \(n\). This wording leaves unresolved whether the intended range is all positive integers \(n\), only primes, or both versions of the question. The displayed assertion defines the question for each positive integer \(n\) without choosing between those ranges. The conclusion preserves the numerical inequality \(\exp(G)\leq n\) exactly as supplied. It has not been replaced by \(\exp(G)\mid n\) or by \(g^n=1_G\) for every \(g\in G\), which are generally different conditions.

For the supplied contextual remarks, let \(F_k\) be the free group on \(k\) generators and define the free Burnside group
\[
B(k,n)=F_k/\langle\!\langle w^n:w\in F_k\rangle\!\rangle,
\]
where the denominator is the normal subgroup generated by all \(n\)-th powers. The source states that, for fixed \(k,n\), a negative answer implies, for example, that \(B(k,n)\) is infinite. It also states that a positive answer, for sufficiently large \(n\), gives an opportunity to find a hyperbolic group \(H\) that is not residually finite and that has a hyperbolic subgroup \(K\) of finite index with no proper subgroup of finite index. The latter subgroup condition means
\[
[H:K]<\infty,\qquad
\forall L\leq K,\quad [K:L]<\infty\Longrightarrow L=K.
\]
A group is residually finite if every nonidentity element has nonidentity image under some homomorphism to a finite group. A hyperbolic group means a finitely generated group with a Cayley graph whose geodesic triangles are uniformly thin. These contextual implications are reported as supplied source claims, not established here; the source gives neither a threshold for ``sufficiently large'' nor precise hypotheses for the stated construction opportunity.
\subsection{Short English statement}
For each fixed number \(k\geq 2\) of generators and each allowed \(n\), must there be a uniform length \(N(k,n)\) such that testing \(n\)-th powers of all products of at most \(N(k,n)\) generators forces every finite group under consideration to have exponent at most \(n\)? The intended restriction of \(n\) to primes and the meaning of the exponent bound require review. The supplied record labels the problem open; that label is not independently verified.
\subsection{Sources}
\begin{itemize}
\item Supplied local UnsolvedMath registry record, permanent ID \(700882\), exact source problem code \texttt{catalogue entry}, title \emph{Kourovka Notebook Problem 12.64}. The supplied status is \texttt{open}; upstream revision: \texttt{d371ee917ac9647fbb06cd6fc529deec6f3dda82}. The record was created and updated on \texttt{2026-10-16T14:00:00Z}.
\item Attributed in the supplied record to A.\ Yu.\ Olshanskii, proposed in 1992: \emph{The Kourovka Notebook (Unsolved Problems in Group Theory)}, original problem 12.64, 12th issue (1992). The record identifies the consulted edition as the 21st issue edition (2026), \texttt{arXiv:1401.0300v48}, and identifies E.\ I.\ Khukhro and V.\ D.\ Mazurov as editors of the notebook collection. This bibliographic attribution is supplied evidence and has not been independently checked.
\item The contextual claims about free Burnside groups and hyperbolic groups are supplied with the problem statement. No separate supporting citation or source-verification date is supplied for those claims.
\end{itemize}
\end{document}
